%
% circuit-runtime.tex
%
% A Z specification of Circuit's runtime admission and session lifecycle.
% The model is intentionally small: it validates the design shape of load,
% scaffold, start, advance, stop, and context injection using totalized
% operation schemas.
%
% Type-check:
%   fuzz -t docs/spec/circuit-runtime.tex
%

\documentclass[a4paper,10pt,fleqn]{article}
\usepackage[margin=1in]{geometry}
\usepackage{fuzz}
\usepackage[colorlinks=true,linkcolor=blue,citecolor=blue,urlcolor=blue]{hyperref}

\begin{document}

\title{Circuit Runtime: A Z Specification}
\author{Punt Labs}
\date{August 2026}
\hypersetup{
  pdftitle={Circuit Runtime: A Z Specification},
  pdfauthor={Punt Labs},
  pdfsubject={Z Specification},
  pdfcreator={fuzz/probcli}
}
\maketitle
\tableofcontents
\newpage

% ============================================================
\section{Scope}
% ============================================================

This specification models Circuit's runtime admission and session lifecycle: a
machine is loaded, its external BOOL facts are validated against check bindings
and the check registry, sessions are started and advanced, and active sessions
are exposed to harness context injection. It is a model of the runtime design,
not of any one workflow machine. Individual workflow machines remain B machines
under \texttt{machines/*.mch}; this document specifies the host runtime that
admits and operates them.

The model follows the totalized-operation style used for command interfaces:
operations do not disappear when preconditions fail. Instead, each public
operation returns a result drawn from a finite result type. Success branches
change state; error or blocked branches preserve state with $\Xi Circuit$.
This mirrors the CLI and tool surfaces, where every command returns either a
successful report or an explicit diagnostic.

\subsection{Non-goals}

The following implementation details are intentionally abstracted away:

\begin{itemize}
  \item Parsing B syntax. A machine is either present in $machines$ or absent.
  \item The internal state graph of each B machine. Advancement is represented
    by an abstract terminal/non-terminal outcome.
  \item Concrete session-id generation. A fresh session identifier is supplied
    as an input and guarded by $s? \notin sessions$.
  \item Shell command execution. Checks are represented by registry entries
    whose result is either stub-false or implemented-true.
  \item File paths and JSON/YAML serialization. The model tracks the facts those
    files encode: bindings, registry entries, and sessions.
\end{itemize}

% ============================================================
\section{Basic Types and Results}
% ============================================================

\begin{zed}
  [MACH, SESSION, CHECK]
\end{zed}

\begin{zed}
  ZBOOL ::= ztrue | zfalse
\end{zed}

\begin{zed}
  CHECKSTATE ::= missing | stubFalse | implementedTrue | badType
\end{zed}

\begin{zed}
  LOADRESULT ::= loadOK | loadUnknownMachine | loadMissingBinding |
    loadMissingRegistry | loadBadReturnType
\end{zed}

\begin{zed}
  SCAFFOLDRESULT ::= scaffoldOK | scaffoldNoop | scaffoldUnknownMachine
\end{zed}

\begin{zed}
  STARTRESULT ::= startOK | startLoadError
\end{zed}

\begin{zed}
  ADVANCERESULT ::= advanceOK | advanceBlocked | advanceNoActiveSession |
    advanceAmbiguousSession | advanceUnknownSession
\end{zed}

\begin{zed}
  STOPRESULT ::= stopOK | stopNoSession | stopAmbiguousSession | stopUnknownSession
\end{zed}

\begin{zed}
  INJECTRESULT ::= injectOK | injectNone
\end{zed}

% ============================================================
\section{State}
% ============================================================

The state schema contains structural invariants: domains line up, active and
terminal sessions are sessions, and the injected context is exactly a set of
sessions. Safety properties are then stated as named consequences of this state
and the totalized operations.

\begin{schema}{Circuit}
  machines : \power MACH \\
  boolFacts : MACH \pfun \power CHECK \\
  bindings : MACH \pfun (CHECK \pfun CHECK) \\
  registry : CHECK \pfun CHECKSTATE \\
  sessions : \power SESSION \\
  active : \power SESSION \\
  terminal : \power SESSION \\
  sessionMachine : SESSION \pfun MACH \\
  injected : \power SESSION
\where
  \dom boolFacts \subseteq machines \\
  \dom bindings \subseteq machines \\
  active \subseteq sessions \\
  terminal \subseteq sessions \\
  active \cap terminal = \emptyset \\
  \dom sessionMachine = sessions \\
  \ran sessionMachine \subseteq machines \\
  injected \subseteq sessions
\end{schema}

\begin{schema}{Init}
  Circuit'
\where
  sessions' = \emptyset \\
  active' = \emptyset \\
  terminal' = \emptyset \\
  sessionMachine' = \emptyset \\
  injected' = \emptyset \\
  ((machines' = \emptyset \\
    \land boolFacts' = \emptyset \\
    \land bindings' = \emptyset \\
    \land registry' = \emptyset) \\
  \lor (machines' = MACH \\
    \land boolFacts' = \{ m : MACH @ m \mapsto CHECK \} \\
    \land bindings' = \emptyset \\
    \land registry' = \emptyset) \\
  \lor (machines' = MACH \\
    \land boolFacts' = \{ m : MACH @ m \mapsto CHECK \} \\
    \land bindings' = \{ m : MACH @ m \mapsto \{ c : CHECK @ c \mapsto c \} \} \\
    \land registry' = \emptyset) \\
  \lor (machines' = MACH \\
    \land boolFacts' = \{ m : MACH @ m \mapsto CHECK \} \\
    \land bindings' = \{ m : MACH @ m \mapsto \{ c : CHECK @ c \mapsto c \} \} \\
    \land registry' = \{ c : CHECK @ c \mapsto implementedTrue \}))
\end{schema}

A machine is loadable when every BOOL fact has a binding, every binding points
at a registry entry, and every referenced registry entry returns BOOL. In this
bounded model, $stubFalse$ and $implementedTrue$ are the two valid BOOL-returning
command cases; $badType$ represents a registry entry with the wrong return type.

\begin{zed}
  BoolRegistry == \{ stubFalse, implementedTrue \}
\end{zed}

\begin{schema}{Loadable}
  Circuit \\
  m? : MACH
\where
  m? \in machines \\
  m? \in \dom boolFacts \\
  m? \in \dom bindings \\
  boolFacts~m? \subseteq \dom (bindings~m?) \\
  \ran (bindings~m?) \subseteq \dom registry \\
  \{ c : \ran (bindings~m?) @ registry~c \} \subseteq BoolRegistry
\end{schema}

% ============================================================
\section{Load}
% ============================================================

\begin{schema}{LoadOK}
  \Xi Circuit \\
  m? : MACH \\
  result! : LOADRESULT
\where
  Loadable \\
  result! = loadOK
\end{schema}

\begin{schema}{LoadUnknownMachine}
  \Xi Circuit \\
  m? : MACH \\
  result! : LOADRESULT
\where
  m? \notin machines \\
  result! = loadUnknownMachine
\end{schema}

\begin{schema}{LoadMissingBinding}
  \Xi Circuit \\
  m? : MACH \\
  result! : LOADRESULT
\where
  m? \in machines \\
  m? \in \dom boolFacts \\
  (m? \notin \dom bindings \lor \lnot (boolFacts~m? \subseteq \dom (bindings~m?))) \\
  result! = loadMissingBinding
\end{schema}

\begin{schema}{LoadMissingRegistry}
  \Xi Circuit \\
  m? : MACH \\
  result! : LOADRESULT
\where
  m? \in machines \\
  m? \in \dom boolFacts \\
  m? \in \dom bindings \\
  boolFacts~m? \subseteq \dom (bindings~m?) \\
  \lnot (\ran (bindings~m?) \subseteq \dom registry) \\
  result! = loadMissingRegistry
\end{schema}

\begin{schema}{LoadBadReturnType}
  \Xi Circuit \\
  m? : MACH \\
  result! : LOADRESULT
\where
  m? \in machines \\
  m? \in \dom boolFacts \\
  m? \in \dom bindings \\
  boolFacts~m? \subseteq \dom (bindings~m?) \\
  \ran (bindings~m?) \subseteq \dom registry \\
  \lnot (\{ c : \ran (bindings~m?) @ registry~c \} \subseteq BoolRegistry) \\
  result! = loadBadReturnType
\end{schema}

\begin{zed}
  LoadMachine \defs LoadOK \lor LoadUnknownMachine \lor LoadMissingBinding
    \lor LoadMissingRegistry \lor LoadBadReturnType
\end{zed}

% ============================================================
\section{Scaffold}
% ============================================================

Scaffold is the operation that makes an incomplete BOOL-fact machine loadable by
creating missing bindings and false stubs. It never overwrites existing registry
entries.


\begin{schema}{ScaffoldOKNewBindings}
  \Delta Circuit \\
  m? : MACH \\
  result! : SCAFFOLDRESULT
\where
  m? \in machines \\
  m? \in \dom boolFacts \\
  boolFacts~m? \neq \emptyset \\
  m? \notin \dom bindings \\
  bindings' = bindings \cup
    \{ m? \mapsto \{ f : boolFacts~m? @ f \mapsto f \} \} \\
  registry' = registry \oplus
    \{ f : boolFacts~m? | f \notin \dom registry @ f \mapsto stubFalse \} \\
  machines' = machines \\
  boolFacts' = boolFacts \\
  sessions' = sessions \\
  active' = active \\
  terminal' = terminal \\
  sessionMachine' = sessionMachine \\
  injected' = injected \\
  result! = scaffoldOK
\end{schema}

\begin{schema}{ScaffoldOKExistingBindings}
  \Delta Circuit \\
  m? : MACH \\
  result! : SCAFFOLDRESULT
\where
  m? \in machines \\
  m? \in \dom boolFacts \\
  boolFacts~m? \neq \emptyset \\
  m? \in \dom bindings \\
  bindings' = bindings \oplus
    \{ m? \mapsto ((bindings~m?) \oplus
      \{ f : boolFacts~m? | f \notin \dom (bindings~m?) @ f \mapsto f \}) \} \\
  registry' = registry \oplus
    \{ f : boolFacts~m? | f \notin \dom registry @ f \mapsto stubFalse \} \\
  machines' = machines \\
  boolFacts' = boolFacts \\
  sessions' = sessions \\
  active' = active \\
  terminal' = terminal \\
  sessionMachine' = sessionMachine \\
  injected' = injected \\
  result! = scaffoldOK
\end{schema}

\begin{schema}{ScaffoldNoop}
  \Xi Circuit \\
  m? : MACH \\
  result! : SCAFFOLDRESULT
\where
  m? \in machines \\
  (m? \notin \dom boolFacts \lor boolFacts~m? = \emptyset) \\
  result! = scaffoldNoop
\end{schema}

\begin{schema}{ScaffoldUnknownMachine}
  \Xi Circuit \\
  m? : MACH \\
  result! : SCAFFOLDRESULT
\where
  m? \notin machines \\
  result! = scaffoldUnknownMachine
\end{schema}

\begin{zed}
  ScaffoldMachine \defs ScaffoldOKNewBindings \lor ScaffoldOKExistingBindings
    \lor ScaffoldNoop \lor ScaffoldUnknownMachine
\end{zed}

% ============================================================
\section{Start}
% ============================================================

Starting a machine is allowed only through the same loadability predicate as the
standalone Load operation. Start therefore cannot create an active session for a
machine with missing bindings, missing registry entries, or non-BOOL registry
entries.

\begin{schema}{StartOK}
  \Delta Circuit \\
  m? : MACH \\
  s? : SESSION \\
  result! : STARTRESULT
\where
  Loadable \\
  s? \notin sessions \\
  sessions' = sessions \cup \{ s? \} \\
  active' = active \cup \{ s? \} \\
  terminal' = terminal \\
  sessionMachine' = sessionMachine \cup \{ s? \mapsto m? \} \\
  machines' = machines \\
  boolFacts' = boolFacts \\
  bindings' = bindings \\
  registry' = registry \\
  injected' = injected \\
  result! = startOK
\end{schema}

\begin{schema}{StartLoadError}
  \Xi Circuit \\
  m? : MACH \\
  s? : SESSION \\
  result! : STARTRESULT
\where
  \lnot (m? \in machines \land m? \in \dom boolFacts \land m? \in \dom bindings
    \land boolFacts~m? \subseteq \dom (bindings~m?)
    \land \ran (bindings~m?) \subseteq \dom registry
    \land \{ c : \ran (bindings~m?) @ registry~c \} \subseteq BoolRegistry) \\
  result! = startLoadError
\end{schema}

\begin{zed}
  StartMachine \defs StartOK \lor StartLoadError
\end{zed}

% ============================================================
\section{Advance}
% ============================================================

Advance is totalized over the command surface. Implicit advance succeeds only
when exactly one session is active. Targeted advance succeeds only for an active
session. A check result of $stubFalse$ blocks the transition; an implemented true
check may advance. Terminal advances remove the session from $active$ and add it
to $terminal$.

\begin{schema}{AdvanceImplicitOK}
  \Delta Circuit \\
  s! : SESSION \\
  terminal? : ZBOOL \\
  result! : ADVANCERESULT
\where
  \# active = 1 \\
  s! \in active \\
  terminal? = zfalse \\
  sessions' = sessions \\
  active' = active \\
  terminal' = terminal \\
  sessionMachine' = sessionMachine \\
  machines' = machines \\
  boolFacts' = boolFacts \\
  bindings' = bindings \\
  registry' = registry \\
  injected' = injected \\
  result! = advanceOK
\end{schema}

\begin{schema}{AdvanceImplicitTerminalOK}
  \Delta Circuit \\
  s! : SESSION \\
  terminal? : ZBOOL \\
  result! : ADVANCERESULT
\where
  \# active = 1 \\
  s! \in active \\
  terminal? = ztrue \\
  sessions' = sessions \\
  active' = active \setminus \{ s! \} \\
  terminal' = terminal \cup \{ s! \} \\
  sessionMachine' = sessionMachine \\
  machines' = machines \\
  boolFacts' = boolFacts \\
  bindings' = bindings \\
  registry' = registry \\
  injected' = injected \setminus \{ s! \} \\
  result! = advanceOK
\end{schema}

\begin{schema}{AdvanceNoActiveSession}
  \Xi Circuit \\
  result! : ADVANCERESULT
\where
  active = \emptyset \\
  result! = advanceNoActiveSession
\end{schema}

\begin{schema}{AdvanceAmbiguousSession}
  \Xi Circuit \\
  result! : ADVANCERESULT
\where
  \# active > 1 \\
  result! = advanceAmbiguousSession
\end{schema}

\begin{zed}
  AdvanceImplicit \defs AdvanceImplicitOK \lor AdvanceImplicitTerminalOK
    \lor AdvanceNoActiveSession \lor AdvanceAmbiguousSession
\end{zed}

\begin{schema}{AdvanceTargetedOK}
  \Delta Circuit \\
  s? : SESSION \\
  terminal? : ZBOOL \\
  result! : ADVANCERESULT
\where
  s? \in active \\
  terminal? = zfalse \\
  sessions' = sessions \\
  active' = active \\
  terminal' = terminal \\
  sessionMachine' = sessionMachine \\
  machines' = machines \\
  boolFacts' = boolFacts \\
  bindings' = bindings \\
  registry' = registry \\
  injected' = injected \\
  result! = advanceOK
\end{schema}

\begin{schema}{AdvanceTargetedTerminalOK}
  \Delta Circuit \\
  s? : SESSION \\
  terminal? : ZBOOL \\
  result! : ADVANCERESULT
\where
  s? \in active \\
  terminal? = ztrue \\
  sessions' = sessions \\
  active' = active \setminus \{ s? \} \\
  terminal' = terminal \cup \{ s? \} \\
  sessionMachine' = sessionMachine \\
  machines' = machines \\
  boolFacts' = boolFacts \\
  bindings' = bindings \\
  registry' = registry \\
  injected' = injected \setminus \{ s? \} \\
  result! = advanceOK
\end{schema}

\begin{schema}{AdvanceTargetedUnknown}
  \Xi Circuit \\
  s? : SESSION \\
  result! : ADVANCERESULT
\where
  s? \notin active \\
  result! = advanceUnknownSession
\end{schema}

\begin{zed}
  AdvanceTargeted \defs AdvanceTargetedOK \lor AdvanceTargetedTerminalOK
    \lor AdvanceTargetedUnknown
\end{zed}

% ============================================================
\section{Stop and Context Injection}
% ============================================================

Stop distinguishes stopped from unloaded. Stopping an active session moves it to
$terminal$, the model's stopped-session set. Stopping a known stopped session is
an idempotent success. Stopping when no active or stopped session exists returns
$stopNoSession$.

\begin{schema}{StopImplicitOK}
  \Delta Circuit \\
  s! : SESSION \\
  result! : STOPRESULT
\where
  \# active = 1 \\
  s! \in active \\
  sessions' = sessions \setminus \{ s! \} \\
  active' = \emptyset \\
  terminal' = terminal \setminus \{ s! \} \\
  sessionMachine' = \{ s! \} \ndres sessionMachine \\
  machines' = machines \\
  boolFacts' = boolFacts \\
  bindings' = bindings \\
  registry' = registry \\
  injected' = injected \setminus \{ s! \} \\
  result! = stopOK
\end{schema}

\begin{schema}{StopStoppedSession}
  \Xi Circuit \\
  s! : SESSION \\
  result! : STOPRESULT
\where
  active = \emptyset \\
  \# terminal = 1 \\
  s! \in terminal \\
  result! = stopOK
\end{schema}

\begin{schema}{StopNoSession}
  \Xi Circuit \\
  result! : STOPRESULT
\where
  active = \emptyset \\
  terminal = \emptyset \\
  result! = stopNoSession
\end{schema}

\begin{schema}{StopAmbiguousSession}
  \Xi Circuit \\
  result! : STOPRESULT
\where
  \# active > 1 \\
  result! = stopAmbiguousSession
\end{schema}

\begin{zed}
  StopImplicit \defs StopImplicitOK \lor StopStoppedSession \lor StopNoSession
    \lor StopAmbiguousSession
\end{zed}

\begin{schema}{StopTargetedOK}
  \Delta Circuit \\
  s? : SESSION \\
  result! : STOPRESULT
\where
  s? \in active \\
  sessions' = sessions \setminus \{ s? \} \\
  active' = active \setminus \{ s? \} \\
  terminal' = terminal \setminus \{ s? \} \\
  sessionMachine' = \{ s? \} \ndres sessionMachine \\
  machines' = machines \\
  boolFacts' = boolFacts \\
  bindings' = bindings \\
  registry' = registry \\
  injected' = injected \setminus \{ s? \} \\
  result! = stopOK
\end{schema}

\begin{schema}{StopTargetedStoppedOK}
  \Xi Circuit \\
  s? : SESSION \\
  result! : STOPRESULT
\where
  s? \in terminal \\
  result! = stopOK
\end{schema}

\begin{schema}{StopTargetedUnknown}
  \Xi Circuit \\
  s? : SESSION \\
  result! : STOPRESULT
\where
  s? \notin active \\
  s? \notin terminal \\
  result! = stopUnknownSession
\end{schema}

\begin{zed}
  StopTargeted \defs StopTargetedOK \lor StopTargetedStoppedOK
    \lor StopTargetedUnknown
\end{zed}

\begin{schema}{InjectActiveSessions}
  \Delta Circuit \\
  result! : INJECTRESULT
\where
  active \neq \emptyset \\
  injected' = active \\
  machines' = machines \\
  boolFacts' = boolFacts \\
  bindings' = bindings \\
  registry' = registry \\
  sessions' = sessions \\
  active' = active \\
  terminal' = terminal \\
  sessionMachine' = sessionMachine \\
  result! = injectOK
\end{schema}

\begin{schema}{InjectNothing}
  \Delta Circuit \\
  result! : INJECTRESULT
\where
  active = \emptyset \\
  injected' = \emptyset \\
  machines' = machines \\
  boolFacts' = boolFacts \\
  bindings' = bindings \\
  registry' = registry \\
  sessions' = sessions \\
  active' = active \\
  terminal' = terminal \\
  sessionMachine' = sessionMachine \\
  result! = injectNone
\end{schema}

\begin{zed}
  InjectContext \defs InjectActiveSessions \lor InjectNothing
\end{zed}

% ============================================================
\section{Safety Properties}
% ============================================================

The central safety invariant is stated separately for readability, but its
conjuncts are consequences the operations are expected to preserve. Keeping it
named lets us use it both as a proof target for the correct design and as a
counter-example target for buggy variants.

\begin{schema}{SafetyOK}
  Circuit
\where
  active \cap terminal = \emptyset \\
  injected \subseteq active \\
  \forall s : active @ sessionMachine~s \in machines \\
  \forall s : active @
    (sessionMachine~s \in \dom boolFacts \land
     sessionMachine~s \in \dom bindings \land
     boolFacts~(sessionMachine~s) \subseteq \dom (bindings~(sessionMachine~s)) \land
     \ran (bindings~(sessionMachine~s)) \subseteq \dom registry \land
     \{ c : \ran (bindings~(sessionMachine~s)) @ registry~c \} \subseteq BoolRegistry)
\end{schema}

\subsection{Buggy variants to validate the model}

A useful model must detect the faults the runtime design is meant to prevent.
The following variants are not part of the correct operation set. They are the
mutations we should run as separate tiny documents or guarded ProB goals:

\begin{enumerate}
  \item \textbf{StartWithoutLoadValidation}: creates an active session for a
    machine even when \texttt{LoadMachine} would return an error. Expected:
    ProB reaches a state violating \texttt{SafetyOK}.
  \item \textbf{AdvanceImplicitChoosesOne}: with two active sessions, advances
    one instead of returning \texttt{advanceAmbiguousSession}. Expected: a goal
    state records an ambiguous implicit advance.
  \item \textbf{AdvanceTerminalKeepsActive}: terminal advance leaves the session
    in \texttt{active}. Expected: \texttt{active \cap terminal \neq \emptyset}
    is reachable.
  \item \textbf{StopTargetedClearsAll}: targeted stop removes every session.
    Expected: a cross-session mutation latch becomes reachable.
  \item \textbf{ScaffoldCreatesTrueStub}: scaffold inserts
    \texttt{implementedTrue} instead of \texttt{stubFalse}. Expected: a
    scaffolded, unimplemented check can advance.
\end{enumerate}

% ============================================================
\section{Verification Plan}
% ============================================================

This spec type-checks through z-spec:

\begin{verbatim}
z-spec check docs/spec/circuit-runtime.tex
\end{verbatim}

The first ProB run follows the mixed-scope style used in other Punt Labs Z
specs. The Makefile names the scope values as `RUNTIME_SPEC_MACHINES`,
`RUNTIME_SPEC_SESSIONS`, and `RUNTIME_SPEC_CHECKS`. The default run uses two
machines, two sessions, and two checks:

\begin{itemize}
  \item two machines exercise multiple machine configurations;
  \item two sessions exercise implicit ambiguity and targeted operations;
  \item two checks exercise multi-fact binding/registry coverage.
\end{itemize}

The spec has multiple intended initial configurations (unknown machine, missing
binding, missing registry, and loadable), so ProB needs
\texttt{MAX\_INITIALISATIONS} above the default to certify the root completely:

\begin{verbatim}
probcli docs/spec/circuit-runtime.tex -model_check \\
  -card MACH 2 -card SESSION 2 -card CHECK 2 \\
  -p MAX_INITIALISATIONS 15
\end{verbatim}

Observed result on 2026-08-09:

\begin{itemize}
  \item all operations covered;
  \item 133 states analysed;
  \item 2189 transitions fired;
  \item no counterexample found;
  \item all open states visited.
\end{itemize}

The Makefile uses \texttt{z-spec check} for type-checking and direct
\texttt{probcli} for the model-checking target because the z-spec CLI does not
yet expose ProB's \texttt{MAX\_INITIALISATIONS} option.

The next validation step is to create separate buggy variants and record the
counterexample traces listed above.

\subsection{Derived test obligations}

\begin{center}
\begin{tabular}{ll}
\textbf{Spec case} & \textbf{Implementation obligation} \\
Load no BOOL facts & \texttt{build-job} loads without check files \\
Load missing binding & \texttt{start review-flow} fails before session creation \\
Scaffold missing facts & checks YAML and false registry stubs are generated \\
False stub & advance is blocked \\
Replace stub & advance is allowed \\
Start valid machine & new active session is created \\
Terminal advance & session auto-stops \\
Two active sessions & implicit advance returns ambiguity \\
Targeted advance & only selected session changes \\
Context injection & all and only active sessions are injected \\
\end{tabular}
\end{center}

\end{document}
